\documentclass[compress, 8pt]{beamer}
% \documentclass[handout, 8pt]{beamer}
\usepackage[utf8]{inputenc}

% Notes to self: \nts[color option]{text that is a note}
\newif\ifshowcomments
% \showcommentstrue % uncomment to show
\input{preambleB}
\usepackage[normalem]{ulem}

%%%%%%%%%%%%%%%%%%%%%%%%%%%%%%%%%%%%%%%%%%%%%%%%%%%%%%%%%%%%%%%%%%%%%%%%%%%%%%
% Paths to figures
\usepackage{subcaption}

% Avoid a warning on mac; can introduce a warning (and maybe a break) on windows
% \usepackage{etex}
% Avoid a warning
\let\Tiny=\tiny

%Information to be included in the title page:
\title{Linear regression + Matching}
\subtitle{ECON726}
\author[]{Tara Sullivan}
\institute{
Tara Sullivan \\
\textit{tara.sullivan55@login.cuny.edu}
}
\date{
\today
\nts{\\
\medskip
Notes/comments are in gray and will be excluded from the presentation.
}
}

% plot graphs with pgfplots
\usepackage{pgfplots}
\pgfplotsset{compat=newest}
\usepgfplotslibrary{groupplots}
\usepgfplotslibrary{dateplot}
\pgfplotsset{compat=newest,
    every axis/.style={
        axis y line*=left,
        axis x line*=bottom,
        % allows for multi-line legend entries
        legend style={cells={align=left}},
        % allows for multi-line titles
        title style={align=center},
    },
}

% For animations using underbrace 
% Source: https://tex.stackexchange.com/a/565866
\usepackage{xparse, letltxmacro}
\LetLtxMacro{\oldunderbrace}{\underbrace}
\DeclareDocumentCommand{\underbrace}{d<> m e{_}}{%
  \IfValueTF{#1}{% IF <overlay-specification> given
    % using global onlside flag, cf. p82 beamer manual v3.59
    \oldunderbrace{#2\onslide<#1>}_{#3}\onslide%
  }{% ELSE
    \oldunderbrace{#2}_{#3}
  }%
}%

\usetikzlibrary{decorations.pathreplacing,calligraphy}

%%%%%%%%%%%%%%%%%%%%%%%%%%%%%%%%%%%%%%%%%%%%%%%%%%%%%%%%%%%%%%%%%%%%%%%%%%%%%%%%
\begin{document}    
%\beamertemplatenavigationsymbolsempty


%%%%%%%%%%%%%%%%%%%%%%%%%%%%%%%%%%%%%%%%%%%%%%%%%%%%%%%%%%%%%%%%%%%%%%%%%%%%%%%%
\begin{frame}
\titlepage
% \tikzset{external/figure name={intro_}}
\end{frame}

%%%%%%%%%%%%%%%%%%%%%%%%%%%%%%%%%%%%%%%%%%%%%%%%%%%%%%%%%%%%%%%%%%%%%%%%%%%%%%%%
\begin{frame}{Housekeeping}

\textbf{Paper selection assignment:}
\begin{itemize}
    \item Great work! Many of you have chosen excellent papers
\end{itemize}

\vspace{3ex}
\textbf{Homework}
\begin{itemize}
    \item Read MHE 3.3 (Matching and Propensity Score)
    \begin{itemize}
        \item Review 3.1. and 3.2., if you're still shaky on the concepts
    \end{itemize}
    \item Read your paper. You won't understand everything; that's okay.
    \begin{itemize}
        \item Maybe try to teach yourself a bit about the methodology. Read the relevant MHE chapter, or start reading about synthetic control
        \item There is a language to papers. You're doing your best to understand it, even if you aren't fluent
    \end{itemize}
    
    \item \textbf{Are you sure you want to do this paper? Does this seem like something you can do?}
\end{itemize}

\vspace{3ex}
\textbf{Today}
\begin{itemize}
    \item Discuss paper selection
    \item Finish linear regression
    \item Introduction to matching
    \item Review examples in jupyter notebooks
    (\href{https://github.com/tara-sullivan/econ726-notebooks/blob/main/notebooks/04-regression.ipynb}{regression}; \href{https://github.com/tara-sullivan/econ726-notebooks/blob/main/notebooks/05-matching.ipynb}{matching})
\end{itemize}


\end{frame}

\begin{frame}{Review - Empirical Project Timeline}

\textbf{Paper selection}: Due \textbf{\sout{September 16} September 22} [5\% of final grade]
\begin{itemize}
\item Done! 
\end{itemize}

\textbf{Project plan}: Due \textbf{October 14} [15\% of final grade] 
\begin{itemize}
\item This paper should be your first effort at addressing points (1) and (2) above. 
Summarize the public policy or business context of the paper. 
Write down the specific question the paper is evaluating.
Discuss the data your are using (both your replication data, and data for your extension)
Discuss the extensions you are considering, and, if necessary, where you will be getting that data. \textbf{New: I'll ask you to verify that you've read the paper}
\end{itemize} 

\textbf{Replication results}: Due \textbf{November 4}: [10\% of final grade]
\begin{itemize}
\item You should submit some of your main replication results. The final results should be formatted as tables and charts in a word document or PDF, submitted on Brightspace. Your code and data should accessible on GitHub.
\end{itemize}

\textbf{Presentation}: Due \textbf{December 2} \& \textbf{December 9} [10\% of final grade]
\begin{itemize}
\item Prepare a 5-10 minute presentation on your paper. Focus on the research question, main result, and one extension. 
\end{itemize}

\textbf{Final Paper}: Due \textbf{December 16} [35\% of final grade]
\begin{itemize}
\item Your paper should contain 5-8 pages of text, and 5-8 pages of figures (double spaced, 12 point font).
You are expected to use python, and must provide the code and data to replicate your findings. 
Your final project should be submitted on Brightspace, and your code and data should be uploaded to github (note: if your data is too large to upload to github, let me know). 
\end{itemize}

\end{frame}

%%%%%%%%%%%%%%%%%%%%%%%%%%%%%%%%%%%%%%%%%%%%%%%%%%%%%%%%%%%%%%%%%%%%%%%%%%%%%%%%
\begin{frame}{Paper selection}

Thoughts on the selection process
\begin{itemize}
    \item How confident are you in your paper? 
    \item Were you able to find a topic you were excited in?
    \item Is there a causal question you'd rather look at, but you couldn't find a paper for?
\end{itemize}

\end{frame}




%%%%%%%%%%%%%%%%%%%%%%%%%%%%%%%%%%%%%%%%%%%%%%%%%%%%%%%%%%%%%%%%%%%%%%%%%%%%%%%
\miniframesoff
\begin{frame}{Outline}
    \tableofcontents
\end{frame}
\miniframeson



%%%%%%%%%%%%%%%%%%%%%%%%%%%%%%%%%%%%%%%%%%%%%%%%%%%%%%%%%%%%%%%%%%%%%%%%%%%%%%%%
\section[Causality]{Regression and causality}
\miniframesoff
\begin{frame}
    \tableofcontents[currentsection]
\end{frame}
\miniframeson
%%%%%%%%%%%%%%%%%%%%%%%%%%%%%%%%%%%%%%%%%%%%%%%%%%%%%%%%%%%%%%%%%%%%%%%%%%%%%%%%


%%%%%%%%%%%%%%%%%%%%%%%%%%%%%%%%%%%%%%%%%%%%%%%%%%%%%%%%%%%%%%%%%%%%%%%%%%%%%%%%
\begin{frame}{Regression and causality}

Regression gives the best linear approximation to the CEF, in a MMSE sense
\begin{itemize}
    \item But when does a regression represent something causal?
\end{itemize}

\vspace{2ex}
A regression is \textbf{causal} when the CEF it approximates is causal
\begin{itemize}
    \item Helpful to think of a CEF being causal by using the potential outcomes framework
    \item CEF is causal when it describes differences in average potential outcomes for a fixed reference population
\end{itemize}

\end{frame}


%%%%%%%%%%%%%%%%%%%%%%%%%%%%%%%%%%%%%%%%%%%%%%%%%%%%%%%%%%%%%%%%%%%%%%%%%%%%%%%%
\begin{frame}{Na{\"i}ve comparison}

Observed difference in earnings equals the average treatment effect (ATE) on the treated plus selection bias:
\begin{align*}
\EE [Y \vert D = 1] - \EE [Y \vert D=0] 
% =& \EE [Y_1 \vert D=1] { - \EE [Y_0 \vert D = 1] + \EE [Y_0 \vert D = 1]}  - \EE [Y_0 \vert D=0] \\
=& \underbrace{\EE [Y_1 - Y_0 \vert D = 1]}_{\text{ATE on treated}} 
+ \underbrace{\EE [Y_0 \vert D = 1] - \EE [Y_0 \vert D=0]}_{\text{Selection bias}}
\end{align*}

Random assignment ($Y_0, Y_1 \perp D$) will solve the selection problem
\begin{itemize}
    \item Causes selection bias to equal zero
\end{itemize}
\begin{align*}
    \EE [Y_1 \vert D = 1] = \EE[Y_1 \vert D = 0] = \EE [Y_1] \quad
    \text{and}& \quad \EE [Y_0 \vert D = 0] = \EE [Y_0 \vert D = 1] = \EE [Y_0] \\
    \implies&
    \EE [Y_0 \vert D = 1] - \EE [Y_0 \vert D = 0] = 0 \\
    \implies&
    \text{Selection bias is 0} \\
    \implies&
    \text{Na{\"i}ve comparison = ATE on treated}
\end{align*}
What do we do when we don't have random assignment?

\end{frame}

%%%%%%%%%%%%%%%%%%%%%%%%%%%%%%%%%%%%%%%%%%%%%%%%%%%%%%%%%%%%%%%%%%%%%%%%%%%%%%%%
\begin{frame}{Conditional Independence Assumption}

Independence of potential outcomes and the treatment:
\begin{itemize}
     \item [$\implies$] ($Y_0, Y_1 \perp D$)
     \item [$\implies$] randomization
\end{itemize} 
What happens when we don't have independence? Meaning the occurrence of $D$ influences the occurrence of $Y_0$ and $Y_1$ (or vice versa)?

\vspace{3ex}
Perhaps we can condition on certain variables $X$, such that the treatment $D_i$ and the potential outcomes $Y_{1i}, Y_{0i}$ LOOK independent
\begin{equation*}
    \{Y_{0i}, Y_{1i}\} \perp D_i \vert X_i
\end{equation*}

\vspace{3ex}

\end{frame}

%%%%%%%%%%%%%%%%%%%%%%%%%%%%%%%%%%%%%%%%%%%%%%%%%%%%%%%%%%%%%%%%%%%%%%%%%%%%%%%%
\begin{frame}{Conditional independence}

Conditional independence assumption: 
\begin{equation*}
    \{Y_{0i}, Y_{1i}\} \perp D_i \vert X_i
\end{equation*}
Intuitively: occurrence of $D_i$ does not influence the occurrence of $Y_{0i}$ or $Y_{1i}$, or vice versa, conditional on knowing $X_i$:
\begin{align*}
    \EE[Y_{0i} \vert D_i, X_i] &= \EE [Y_{0i} \vert X_i] \\
    \EE[Y_{1i} \vert D_i, X_i] &= \EE [Y_{1i} \vert X_i] 
\end{align*}

\begin{columns}
\column{0.4\linewidth}
\begin{figure}
\begin{tikzpicture}
\node [draw, circle] (D) {D};
\node [draw, circle, below =of D] (Y) {Y};
\node [draw, circle, right =of D] (X) {X};
\draw (D) edge[->] (Y);
\draw (X) edge[->] (D);
\draw (X) edge[->] (Y);
\end{tikzpicture}
\end{figure}

\column{0.6\linewidth}
\begin{figure}
\begin{tikzpicture}
\node [draw] (D) {College attendance};
\node [draw, below =of D] (Y) {Earnings};
\node [draw, right =of D] (X) {Ability};
\draw (D) edge[->] (Y);
\draw (X) edge[->] (D);
\draw (X) edge[->] (Y);
\end{tikzpicture}
\end{figure}

\end{columns}

\end{frame}


%%%%%%%%%%%%%%%%%%%%%%%%%%%%%%%%%%%%%%%%%%%%%%%%%%%%%%%%%%%%%%%%%%%%%%%%%%%%%%%%
\begin{frame}{Selection bias under CIA}

Under conditional independence assumption, selection bias disappears conditional on $X_i$:
\begin{align*}
    \EE [Y_{0i} \vert X_i, D_i = 1] - \EE [Y_{0i} \vert X_i, D_i = 0]
    =& \EE [Y_{0i} \vert X_i, D_i = 0] - \EE [Y_{0i} \vert X_i, D_i = 0] \\
    =& 0
\end{align*}
So we can write the difference in observed outcomes \textbf{conditional on $X_i$} as:
\begin{align*}
    \EE [Y_i \vert X_i, D_i = 1] - \EE [Y_i \vert X_i, D_i = 0] 
    &= \EE [Y_{1i} \vert X_i, D_i = 1] - \EE [Y_{0i} \vert X_i, D_i = 0] \\
    &= \EE [Y_{1i} \vert X_i] - \EE [Y_{0i} \vert X_i] \\
    &= \EE [Y_{1i} - Y_{0i} \vert X_i]
\end{align*}
Given the CIA, comparisons of average earnings (conditional on $X_i$) have a causal interpretation
\end{frame}

%%%%%%%%%%%%%%%%%%%%%%%%%%%%%%%%%%%%%%%%%%%%%%%%%%%%%%%%%%%%%%%%%%%%%%%%%
%%%%%%%%%%%%%%%%%%%%%%%%%%%%%%%%%%%%%%%%%%%%%%%%%%%%%%%%%%%%%%%%%%%%%%%%%%%%%%%%
\begin{frame}{When treatment takes on more than two values}

Expanding CIA assumption to relations that can take on more than two values, such as when $s_i = \{0, 1, \dots, 16\}$ equals years of schooling:
\begin{equation*}
    Y_{si} = f_i (s)
\end{equation*}
Individual specific function:
\begin{itemize}
    \item Potential earnings that person $i$ would receive after obtaining $s$ years of education
    \item What would $i$ learn for any value of schooling
\end{itemize}

\vspace{2ex}
If $s$ could only take on two values, $s=12$ and $s=16$, then we're considering the value of college vs. no college:
\begin{equation*}
    Y_{0i} = f_i(12), \quad \quad Y_{1i} = f_i(16)
\end{equation*}

\end{frame}

%%%%%%%%%%%%%%%%%%%%%%%%%%%%%%%%%%%%%%%%%%%%%%%%%%%%%%%%%%%%%%%%%%%%%%%%%%%%%%%%
\begin{frame}{CIA with non-binary treatment}

CIA in this more general set-up:
\begin{equation*}
    Y_{si} \perp s_i \vert X_i, \quad \forall s_i
\end{equation*}
\begin{itemize}
    \item In randomized experiments (RCTs), $s_i$ is randomly assigned based on covariates $X_i$
    \item In observational studies, the CIA means $s_i$ is \textbf{as good as randomly assigned} conditional on $X_i$
\end{itemize}
Under conditional independence assumption, comparisons of average earnings across schooling levels have a causal interpretation:
\begin{equation*}
    \EE [Y_i \vert X_i, s_i = s] - \EE[Y_i \vert X_i, s_i = s - 1] = \EE [f_i(s) - f_i(s - 1) \vert X_i]
\end{equation*}
For example, average causal effect of high school graduation:
\begin{align*}
    \EE [Y_i \vert X_i, s_i = 12] - \EE[Y_i \vert X_i, s_i = 11] =& \EE [f_i (12) \vert X_i, s_i = 12] - \EE [f_i (11) \vert X_i, s_i = 11] \\
    =& \EE [f_i(12) - f_i(11) \vert X_i, s_i = 12] 
\end{align*}
\begin{itemize}
    \item Here, selection bias would come from the differences in potential earnings of high school graduates and non-graduates
    \item Under CIA, high school graduation is independent of potential earnings conditional on $X_i$, so selection bias disappears
\end{itemize}

\end{frame}


% %%%%%%%%%%%%%%%%%%%%%%%%%%%%%%%%%%%%%%%%%%%%%%%%%%%%%%%%%%%%%%%%%%%%%%%%%%%%%%%%
% \begin{frame}{}

% Unconditional average effects under CIA:
% \begin{align*}
%     \EE [Y_i \vert X_i, s_i = 12] - \EE[Y_i \vert X_i, s_i = 11] =& \EE [f_i (12) \vert X_i, s_i = 12] - \EE [f_i (11) \vert X_i, s_i = 11] \\
%     =& \EE [f_i(12) - f_i(11) \vert X_i, s_i = 12]
% \end{align*}
% Computationally, what does this mean?
% \begin{itemize}
%     \item 
% \end{itemize}

% \end{frame}


%%%%%%%%%%%%%%%%%%%%%%%%%%%%%%%%%%%%%%%%%%%%%%%%%%%%%%%%%%%%%%%%%%%%%%%%%%%%%%%%
\begin{frame}{CIA and regression}

Regression is an empirical strategy that turns CIA into causal effects:
\begin{itemize}
    \item Assume $f_i(s)$ is linear in $s$, and the same for everyone, except for an additive term $\eta_i$:
    \begin{equation*}
        f_i(s) = \alpha + \rho s + \eta_i
    \end{equation*}
    \item \textbf{Linear constant effects model}: equation is linear, and the functional relationship is the same for everyone
    \item Only individual-specific part of $f_i(s)$ is mean-zero $\eta_i$, which captures unobserved factors that determine potential earnings
\end{itemize}
Observed value will be:
\begin{equation*}
    Y_i = \alpha + \rho s_i + \eta_i, 
\end{equation*}
Note that $s_i$ is correlated with $\eta_i$: $\EE [\eta_i \vert s_i] \neq 0$. If we ran this regression, we'd expect selection bias to impact our results.

\end{frame}

%%%%%%%%%%%%%%%%%%%%%%%%%%%%%%%%%%%%%%%%%%%%%%%%%%%%%%%%%%%%%%%%%%%%%%%%%%%%%%%%
\begin{frame}{CIA and regression (cont.)}

Assume CIA holds, and $\eta_i$ has the form:
\begin{equation*}
    \eta_i = X_i^\prime \gamma + v_i
\end{equation*}
where: 
\begin{itemize}
    \item $\gamma$ is a vector of population regression coefficients that satisfies $\EE [\eta_i \vert X_i] = X_i^\prime \gamma$
    \item $v_i$ and $X_i$ are uncorrelated by construction
\end{itemize}
Then, because of the CIA:
\begin{align*}
    \EE [f_i(s) \vert X_i, s_i] &= \EE[f_i (s) \vert X_i] \\
    &= \alpha + \rho s + \EE [\eta_i \vert X] \\
    &= \alpha + \rho s + X_i^\prime \gamma
\end{align*}
The residual in the linear causal model $v_i$ is uncorrelated with $s_i$ and $X_i$, and $\rho$ is the causal effect of interest:
\begin{equation*}
    Y_i = \alpha + \rho s_i + X_i^\prime \gamma + v_i
\end{equation*}
$\implies$ if we assume CIA holds in this linear causal model, $\rho$ is a causal coefficient
\end{frame}


%%%%%%%%%%%%%%%%%%%%%%%%%%%%%%%%%%%%%%%%%%%%%%%%%%%%%%%%%%%%%%%%%%%%%%%%%%%%%%%%
\begin{frame}{Regression takeaways}

We just discussed how CIA + linear causal model yields a causal interpretation of CEF
\begin{itemize}
    \item Linear causal model: $Y_i = \alpha + \rho s_i + \eta_i$
    \item If CIA holds given some variables $X_i$: $Y_{0i}, Y_{1i} \perp s_i \vert X_i$
    \item Our regression function has a causal interpretation: $Y_i = \alpha + \rho s_i + X_i^\prime \gamma + v_i$
    \begin{itemize}
        \item $\rho$ is the causal effect of $s_i$ on $Y_i$
    \end{itemize}
    \item CEF given by:
    \begin{equation*}
        \EE [Y_i \vert X_i, s_i] = \alpha + \rho s_i + X_i^\prime \gamma
    \end{equation*}
\end{itemize}

% \textbf{TO DO:} regression example in pyhton; use companion website
\vspace{3ex}
\textbf{Thing to note:}
\begin{itemize}
    \item CIA is the same as unconfoundedness is the same as selection on observables
    \item In regression, CIA is the same as exogeneity
\end{itemize}

\end{frame}



%%%%%%%%%%%%%%%%%%%%%%%%%%%%%%%%%%%%%%%%%%%%%%%%%%%%%%%%%%%%%%%%%%%%%%%%%%%%%%%%
\section[Matching]{Matching}
\miniframesoff
\begin{frame}
    \tableofcontents[currentsection]
\end{frame}
\miniframeson
%%%%%%%%%%%%%%%%%%%%%%%%%%%%%%%%%%%%%%%%%%%%%%%%%%%%%%%%%%%%%%%%%%%%%%%%%%%%%%%%

% %%%%%%%%%%%%%%%%%%%%%%%%%%%%%%%%%%%%%%%%%%%%%%%%%%%%%%%%%%%%%%%%%%%%%%%%%%%%%%%%
% \begin{frame}{}



% Linear causal model seems like a strong assumption. How realistic is it to assume that policy responses are linear?
% \begin{itemize}
%     \item Luckily, you don't always need to worry about that to run regressions
%     \item You can think of linear regression of $Y_i$ on $X_i$ and $s_i$ as providing the best \textbf{linear approximation} to the underlying CEF
%     \item 
% \end{itemize}

% \end{frame}

%%%%%%%%%%%%%%%%%%%%%%%%%%%%%%%%%%%%%%%%%%%%%%%%%%%%%%%%%%%%%%%%%%%%%%%%%%%%%%%%
\begin{frame}{Motivating example}

Consider the universe of applicants to the U.S. military:
\begin{itemize}
    \item $Y_i$: average earnings
    \item $D_i$: dummy variable indicating whether $i$ served in the military
\end{itemize}
We may be interested in the average effect of treatment on the treated (ToT):
\begin{equation*}
    \EE [Y_{1i} - Y_{0i]} \vert D_i = 1]
\end{equation*}
Difference between:
\begin{itemize}
    \item Average earnings of soldiers (observable)
    \item Average earnings of soldiers, had they not served (unobservable)
\end{itemize}
Na{\"i}ve comparisons will be biased:
\begin{align*}
\EE [Y_i \vert D_i = 1] - \EE [Y_i \vert D_i=0] 
=& \underbrace{\EE [Y_{1i} - Y_{0i} \vert D_i = 1]}_{\text{ATE on treated}} 
+ \underbrace{\EE [Y_{0i} \vert D_i = 1] - \EE [Y_{0i} \vert D_i=0]}_{\text{Selection bias}}
\end{align*}

\end{frame}

%%%%%%%%%%%%%%%%%%%%%%%%%%%%%%%%%%%%%%%%%%%%%%%%%%%%%%%%%%%%%%%%%%%%%%%%%%%%%%%%
\begin{frame}{Motivating example: CIA}

Conditional independence assumption:
\begin{equation*}
    Y_{0i}, Y_{1i} \perp D_i \vert X_i
\end{equation*}
Assuming CIA, selection bias disappears after conditioning on $X_i$:
\begin{align*}
    \delta_{ToT} &= \EE[Y_{1i} - Y_{0i} \vert D_i = 1] \\
    \onslide<2->{
    &= \EE \left[ \left. 
    \EE[Y_{1i} - Y_{0i} \alert<2>{\vert X_i,} D_i = 1]
    \right\vert D_i = 1
     \right] &\because \text{LIE}
     } \\
    \onslide<3->{
        &= \EE \left[ \left.
            \EE [Y_{1i} \vert X_i, D_i=1]
            - \alert<4>{\EE [Y_{0i} \vert X_i, D_i=1]}
            \right\vert D_i = 1
        \right]
    } \\
    \onslide<4->{
        &= \EE \left[ \left.
            {\color<5>{Green}\EE [Y_{1i} \vert X_i, D_i=1]}
            - \alert<4>{{\color<5>{Orange} \EE [Y_{0i} \vert X_i, D_i=0]}}
            \right\vert D_i = 1
        \right]        
    } \\
    \onslide<6->{
        &= \EE \left[ \left.
            {\color<6>{Green}\EE [Y_{i} \vert X_i, D_i=1]}
            - {\color<6>{Orange} \EE [Y_{i} \vert X_i, D_i=0]}              
            \right\vert D_i = 1
        \right]
    }
\end{align*}
\onslide<4->{
    By CIA:
    \begin{equation*}
        \alert<4>{\EE [Y_{0i} \vert X_i, D_i = 0]
        = \EE [Y_{0i} \vert X_i, D_i = 1]}
    \end{equation*}
}
\onslide<5->{
    By definition of $Y_{i} = Y_{0i} + D_i (Y_{1i} - Y_{0i})$:
    \begin{align*}
        {\color<6>{Orange}\EE [Y_{i} \vert X_i, D_i = 0]} 
        &= \EE [Y_{0i} + D_i (Y_{1i} - Y_{0i}) \vert X_i, D_i = 0] \\
        &= {\color<5-6>{Orange} \EE [Y_{0i} \vert X_i, D_i = 0] } \\
        {\color<6>{Green}\EE [Y_{i} \vert X_i, D_i = 1]} 
        &= {\color<5-6>{Green} \EE [Y_{1i} \vert X_i, D_i = 1]} 
    \end{align*}
}
\end{frame}

%%%%%%%%%%%%%%%%%%%%%%%%%%%%%%%%%%%%%%%%%%%%%%%%%%%%%%%%%%%%%%%%%%%%%%%%%%%%%%%%
\begin{frame}{}

\begin{align*}
    \delta_{ToT}
        &= 
        \EE \left[ \left.
            \underbrace<4>{
            \underbrace<2>{\color<2>{Green}\EE [Y_{i} \vert X_i, D_i=1]}_{}
            - \underbrace<3>{\color<3>{Orange}\EE [Y_{i} \vert X_i, D_i=0]}_{}
            }_{}
            \right\vert D_i = 1
        \right] \\
        \onslide<5->{
        &= \EE \Big[ {\color<5>{MainBlue}\delta_X} \Big\vert D_i=1\Big]
        } 
\end{align*}
Consider the parts of this estimator:
\begin{itemize}
    \item[]<2-> {\color<2>{Green}$\EE [Y_{i} \vert X_i, D_i=1]$}: \quad \textbf{For a given $X_i$}, take the average over all observed \textbf{treated} $Y_i$

    \item[]<3-> {\color<3>{Orange}$\EE [Y_{i} \vert X_i, D_i=0]$}: \quad \textbf{For a given $X_i$}, take the average over all observed \textbf{untreated} $Y_i$
\end{itemize}
\onslide<4->{
\vspace{2ex}
    Define $\delta_X$ to be the average difference in oberved $Y_i$ for treated and untreated with a given value $X_i$:
    \begin{equation*}
        {\color<5>{MainBlue}\delta_X} \equiv
        \underbrace<4>{\EE [Y_{i} \vert X_i, D_i=1]
        - \EE [Y_{i} \vert X_i, D_i=0]}_{}
    \end{equation*}
    \begin{itemize}
        \item Example: difference in mean earnings by veteran status at each value of $X_i$
    \end{itemize}
}
% \onslide<5->{
%     Effect of treatment on the treated can be written as:
%     \begin{equation*}
%         \delta_ToT = \EE \left[ . \right\vert D_i \right]
%     \end{equation*}
% }

\end{frame}

%%%%%%%%%%%%%%%%%%%%%%%%%%%%%%%%%%%%%%%%%%%%%%%%%%%%%%%%%%%%%%%%%%%%%%%%%%%%%%%%
\begin{frame}{}

Effect of treatment on the treated under CIA: 
\begin{align*}
    \delta_{ToT} &= \EE \left[ \left.
            {\EE [Y_{i} \vert X_i, D_i=1]}
            - {\EE [Y_{i} \vert X_i, D_i=0]}        
            \right\vert D_i = 1
        \right] \\
        &= \EE \left[ \left. 
            \delta_X
        \right\vert D_i = 1\right] \\
    \text{where} & \\
        \delta_X &= {\EE [Y_{i} \vert X_i, D_i=1]}
            - {\EE [Y_{i} \vert X_i, D_i=0]}
\end{align*}
Let's say there was some variable $X_i$ that represents an indicator of whether an applicant is ``high ability'' or ``lower ability''
\begin{itemize}
    \item $X_i = 0$ means $i$ is low ability 
    \item $X_i = 1$ means $i$ is high ability
\end{itemize}
We'll assume CIA: $\{Y_{0i}, Y_{1i}\} \perp D_i \vert X_i$ 
\begin{itemize}
    \item Intuitively, what does this mean? Is this a good assumption? Is it realistic?
\end{itemize}

The matching estimand can be written as:
\begin{align*}
    \delta_{ToT} = \EE [Y_{1i} - Y_{0i} \vert D_i = 1] &= \EE [\delta_X \vert D_i = 1] \\
    &= \sum_{x \in \{0, 1\}} \delta_x \PP(X_i = x \vert D_i = 1 ) \\
    &= \delta_0 \PP (X_i = 0 \vert D_i = 1)
    + \delta_1 \PP (X_i = 1 \vert D_i = 1)
\end{align*}
\end{frame}

%%%%%%%%%%%%%%%%%%%%%%%%%%%%%%%%%%%%%%%%%%%%%%%%%%%%%%%%%%%%%%%%%%%%%%%%%%%%%%%%
\begin{frame}{}

\begin{align*}
    \delta_{ToT} &= \EE [\delta_X \vert D_i = 1] \\
    &= 
        {\color<2>{Orange}\delta_{X_i = 0}} 
        {\color<4>{Orange}\PP (X_i = 0 \vert D_i = 1)}
    + 
        {\color<3>{Green}\delta_{X_i = 1}} 
        {\color<5>{Green}\PP (X_i = 1 \vert D_i = 1)}
\end{align*}
How do these elements connect to data? 
\begin{table}
\begin{tabular}{rl}
\onslide<2->{
    {\color<2>{Orange} $\delta_{X_i = 0}$}
    &Difference in observed $Y_i$ for treated and untreated with $X_i = 0$\\
    &{\small\color{DarkGray}(Average difference in wages by veteran status for \textbf{low} ability $i$)} \\
}
\onslide<3->{
    {\color<3>{Green} $\delta_{X_i = 1}$}
    &Difference in observed $Y_i$ for treated and untreated with $X_i = 1$\\
    &{\small\color{DarkGray}(Average difference in wages by veteran status for \textbf{high} ability $i$)} \\
}
\onslide<4->{
    {\color<4>{Orange} $\PP (X_i = 0 \vert D_i = 1)$}
    &Probability $X_i = 0$ conditional on being treated\\
    &{\small\color{DarkGray}(Fraction of veterans that are \textbf{low} ability)} \\
}
\onslide<5->{
    {\color<5>{Green} $\PP (X_i = 1 \vert D_i = 1)$}
    &Probability $X_i = 1$ conditional on being treated\\
    &{\small\color{DarkGray}(Fraction of veterans that are \textbf{high} ability)} \\
}

\end{tabular}
\end{table}
\onslide<6->
{Intuitively, what is this matching estimate?
\begin{itemize}
    \item $X_i$-specific treatment-control comparisons, weighted together to produce a single overall treatment effect
\end{itemize}
}
\end{frame}

%%%%%%%%%%%%%%%%%%%%%%%%%%%%%%%%%%%%%%%%%%%%%%%%%%%%%%%%%%%%%%%%%%%%%%%%%%%%%%%%
\begin{frame}{}

Effect of treatment on the treated under CIA: 
\begin{align*}
    \delta_{ToT} &= \EE \left[ \left.
            {\EE [Y_{i} \vert X_i, D_i=1]}
            - {\EE [Y_{i} \vert X_i, D_i=0]}        
            \right\vert D_i = 1
        \right] \\
        &= \EE \left[ \left. 
            \delta_X
        \right\vert D_i = 1\right] \\
    \text{where} & \\
        \delta_X &= {\EE [Y_{i} \vert X_i, D_i=1]}
            - {\EE [Y_{i} \vert X_i, D_i=0]}
\end{align*}
Now instead of $X_i$ being a dummy, assume $X_i$ is the result of test of ability/intelligence
\begin{itemize}
    \item $X_i \in \{0, 1, \dots, 100\}$ and measures ability
\end{itemize}
We'll assume CIA: $\{Y_{0i}, Y_{1i}\} \perp D_i \vert X_i$ 
\begin{itemize}
    \item Intuitively, what does this mean? Is this a good assumption? Is it realistic?
\end{itemize}
The matching estimand can be written as:
\begin{align*}
    \delta_{ToT} %= \EE [Y_{1i} - Y_{0i} \vert D_i = 1] 
    &= \EE [\delta_X \vert D_i = 1] \\
    &= \sum_{x=0}^{100} \delta_x \PP(X_i = x \vert D_i = 1 ) \\
    &= \delta_0 \PP (X_i = 0 \vert D_i = 1)
    + \delta_1 \PP (X_i = 1 \vert D_i = 1)
    + \dots + \delta_{100} \PP (X_i = 100 \vert D_i = 1)
\end{align*}
Intuitively, what is this matching estimate?
\begin{itemize}
    \item $X_i$-specific treatment-control comparisons, weighted together to produce a single overall treatment effect
\end{itemize}

\end{frame}

%%%%%%%%%%%%%%%%%%%%%%%%%%%%%%%%%%%%%%%%%%%%%%%%%%%%%%%%%%%%%%%%%%%%%%%%%%%%%%%%
\begin{frame}{}

Now instead of $X_i$ isn't a single variable, but instead a vector of (discrete) variables
\begin{itemize}
    \item Year of birth, test score, year of application to the military, educational attainment at the time of application
\end{itemize}
Think of each realized $X_i=x$ as a particular combination of all these variables
\begin{itemize}
    \item {\color{DarkGray}For example, $X_i = z$ might describe an applicant born in 1970, who scored a 95 on an ability test, and graduated from college}
\end{itemize}
We'll assume CIA: $\{Y_{0i}, Y_{1i}\} \perp D_i \vert X_i$ 
\begin{itemize}
    \item Intuitively, what does this mean? Is this a good assumption? Is it realistic?
\end{itemize}
The matching estimand can be written as:
\begin{align*}
    \delta_{ToT} %= \EE [Y_{1i} - Y_{0i} \vert D_i = 1] 
    &= \EE [\delta_X \vert D_i = 1] \\
    &= \sum_{x} \delta_x \PP(X_i = x \vert D_i = 1 ) 
\end{align*}
Intuitively, what is this matching estimate?
\begin{itemize}
    \item The $\delta_{x}$ are $x$-specific treatment-control comparisons
    \begin{itemize}
        \item {\color{DarkGray}For $X_i = z$, compare average wage difference of veterans and non-veterans who were born in 1970, who scored a 95 on an ability test, and graduated from college}
        \item Difference in average wages by veteran status for those with the same $X_i$
    \end{itemize}
    \item Weighted together to produce a single overall treatment effect
    \begin{itemize}
        \item {\color{DarkGray} Calculate weights based on the distribution of characteristics $X_i$ among veterans}
    \end{itemize}
\end{itemize}
% Note: this basically generalizes to continuous $X_i$, but the math is much harder

\end{frame}

%%%%%%%%%%%%%%%%%%%%%%%%%%%%%%%%%%%%%%%%%%%%%%%%%%%%%%%%%%%%%%%%%%%%%%%%%%%%%%%%
\begin{frame}{Unconditional Average Treatment Effect}

Conditional treatment effect on treated:
\begin{align*}
    \delta_{ToT} = \EE [Y_{1i} - Y_{0i} \vert D_i = 1] &= \EE [\delta_X \vert D_i = 1] \\
    &= \sum_{x \in \{0, 1\}} \delta_x \PP(X_i = x \vert D_i = 1 ) \\
    &= \delta_0 \PP (X_i = 0 \vert D_i = 1)
    + \delta_1 \PP (X_i = 1 \vert D_i = 1)
\end{align*}
\begin{itemize}
    \item How much the typical soldier gained or lost as a result of enlisting
\end{itemize}

\vspace{2ex}
Can also write unconditional average treatment effect:
\begin{align*}
    \delta_{ATE} &= \EE \left[ 
        \EE [Y_{1i} \vert X_i, D_i = 1]
        - \EE [Y_{0i} \vert X_i, D_i = 0]
    \right] \\
    &= \EE [Y_{1i} - Y_{0i}] \\
    &= \sum_x \delta_x \PP(X_i = x)
\end{align*}
\begin{itemize}
    \item How much the typical applicant would gain or lose as a result of enlisting
\end{itemize}

\end{frame}





%%%%%%%%%%%%%%%%%%%%%%%%%%%%%%%%%%%%%%%%%%%%%%%%%%%%%%%%%%%%%%%%%%%%%%%%%%%%%%%%
\begin{frame}{Matching and regression}

Matching and regression are related:
\begin{itemize}
    \item Both are comparing outcomes (differences in $Y_i$) over treatment groups, controlling for covariates $X_i$
    \item These methods differ in implementation
\end{itemize}
Regression estimates of matching strategy:
\begin{equation*}
    Y_i = \sum_x d_{ix} \alpha_x + \delta_R D_i + e_i
\end{equation*}
where:
\begin{itemize}
    \item $d_{ix} = \mathbf{1} [X_i = x]$: dummy variable that equals 1 if $X_i = x$, 0 otherwise
    \item $\alpha_x$: regression effect for $X_i = x$
    \item $\delta_R$: regression estimand
\end{itemize}

\end{frame}

%%%%%%%%%%%%%%%%%%%%%%%%%%%%%%%%%%%%%%%%%%%%%%%%%%%%%%%%%%%%%%%%%%%%%%%%%%%%%%%%
\begin{frame}{Matching and regression differences}

In general, estimates from matching and from regression will often be similar, but not identical.
\begin{itemize}
    \item Regression can be thought of as a type of matching strategy
    \begin{itemize}
        \item Average effects are calculated by covariates $X$: $\delta_X$
        \item These effects are weighed to produce a weighted average
    \end{itemize}

    \item Difference in how weights are determined. For ATE on treated:
    \begin{itemize}
        \item Matching: weights are determined by distribution of covariates on treated
        \item Regression: uses treatment-variance weighted average
        \item These will differ unless treatment is independent of covariates
    \end{itemize}

\end{itemize}

\vspace{3ex}
Note: neither regression nor covariate matching estimands give any weights to covariate cells that do not contain treated and control observations
\begin{itemize}
    \item Often called \textbf{common support}
\end{itemize}
\end{frame}





\end{document}
